\begin{definition}\label{long}\textbf{Longitud de curva.}
    Sean $C$ una curva simple en $\Rn{n}$ y  $\boldsymbol{\sigma}:I\subseteq\R\to\Rn{n}$ una parametrizaci\'on regular y  de clase $\mathcal{C}^1$ a trozos de $C$  (ver\: \ref{param}). Se define la \emph{longitud de} $C$, y se nota  Long($C$), a 
    \[
        \text{Long}(C)=\sum_{i=1}^{m}\int_{I_i}\|\boldsymbol{\sigma}'(t)\|dt,    
    \]
   donde $ I_i= (t_i, t_{i+1})$,  $\boldsymbol{\sigma}\lvert_{I_i}$ es de clase $\mathcal{C}^1$  y  $\{t_i\}_{i=0}^{m-1}$  es una partici\'on de $I$.
\end{definition}



\begin{example}
    Sea la curva simple $\Gamma$ que describe una circunferencia de radio $r$ centrada en el origen, \emph{i.e.}
        \[
            \Gamma = \bigl\{(x,y)\in\Rn{2} \;:\; x^2 + y^2 = r^2 \bigr\}.
        \]
         Calcular la longitud de $\Gamma$.
        
        
       Es sencillo ver que 
        \[
            \boldsymbol{\sigma}(t) 
            = 
            \bigl(r\cos t,\;r\sin t\bigr),
            \qquad t\in [\,0,\,2\pi\,].
        \] 
       es una parametrización regular y de clase  \(\mathcal{C}^1\)  de \(\Gamma\), luego  
        \begin{gather*}
            \boldsymbol{\sigma}'(t)
            =
            \bigl(-r\sin t,\;r\cos t\bigr) \mbox{ y }\\[.2cm]
            \|\boldsymbol{\sigma}'(t)\|
            =
            \sqrt{(-r\sin t)^2 + (r\cos t)^2}
            =
            r\,\sqrt{\sin^2 t + \cos^2 t}
            = 
            r.
        \end{gather*}
       Para calcular la longitud de $\Gamma$  basta con plantear y resolver la siguiente integral
        \[
            \mathrm{Long}(\Gamma)
            =
            \int_{0}^{2\pi} \bigl\|\boldsymbol{\sigma}'(t)\bigr\| \, dt
            =
            \int_{0}^{2\pi} r \; dt
            =
            2\pi r.
        \]
        Por lo tanto, la  longitud de la circunferencia de radio \(r\) es \(2\pi r\), resultado ya conocido. 
    \end{example}
    
   
    \begin{example}
    Sea  la curva simple $C$ dada por el segmento que une los puntos $\mathbf{p}$ y $\mathbf{q}$  en \(\Rn{n}\).   Calcular la longitud de $C$. 
    
     Es sencillo ver que 
        \[
            \boldsymbol{\sigma}(t) 
            = 
            \mathbf{p} + t\, (\,\mathbf{q} - \mathbf{p}\,),
            \qquad t\in [\,0,\,1\,].
        \]
          es una parametrización regular y de clase  \(\mathcal{C}^1\) de \(C\), luego  
        \[
            \boldsymbol{\sigma}'(t) = \mathbf{q} - \mathbf{p}
            \implies \|\boldsymbol{\sigma}'(t)\| = \|\mathbf{q} - \mathbf{p}\|
        \]
      Para calcular la longitud de $\Gamma$  basta con plantear y resolver la siguiente integral  
        \[
            \mathrm{Long}(C)
            =
            \int_{0}^{1} \|\boldsymbol{\sigma}'(t)\|\: dt
            =
            \int_{0}^{1} \|\mathbf{q} - \mathbf{p}\|\: dt
            =
            \|\mathbf{q} - \mathbf{p}\|.
        \] Por lo tanto, la longuitud del segmento que une los puntos  $\mathbf{p}$ y $\mathbf{q}$  es $ \|\mathbf{q} - \mathbf{p}\|$, resultado ya conocido.
 \end{example}

 \begin{tarea} Calcular, usando la definici\'on  \ref{long},   el per\'imetro de un rombo cuyas diagonales menor y mayor miden $d$ y $D$ respectivamente.  
 \end{tarea}


\begin{definition}\textbf{Integral de trayectoria de campos escalares.}
    Sea $\boldsymbol{\sigma}:I=[a,b]\to\Rn{n}$ una trayectoria de clase $\mathcal{C}^1(\interior{I})$ y sea $f$ un campo escalar tal que este bien definida y sea continua   la funci\'on $f\circ\boldsymbol{\sigma}$.  Se define la \emph{integral de trayectoria}, o \emph{integral de} $f$ \emph{a lo largo de la trayectoria} $\boldsymbol{\sigma}$, y se la nota $\int_{\boldsymbol{\sigma}}f\:ds$, a
    \[
        \int_{\boldsymbol{\sigma}}f\:ds=\int_a^b (f\circ\boldsymbol{\sigma})(t)  \|\boldsymbol{\sigma}'(t)\|\:dt.
    \]
\end{definition}

Si $\boldsymbol{\sigma}(t)$ es de clase $\mathcal{C}^1$ a trozos, se define $\int_{\boldsymbol{\sigma}}f\:ds$  de la siguiente manera:  
\begin{definition}   Sea $\boldsymbol{\sigma}:I=[a,b]\to\Rn{n}$ una trayectoria de clase $\mathcal{C}^1$ a trozos   y  sea $f$ un campo escalar tal que este bien definida y sea continua   la funci\'on $f\circ\boldsymbol{\sigma}$.  Se define la \emph{integral de trayectoria}  o \emph{integral de} $f$ \emph{a lo largo de la trayectoria} $\boldsymbol{\sigma}$, y se  nota $\int_{\boldsymbol{\sigma}}f\:ds$, a
  
    \[
       \int_{\boldsymbol{\sigma}}f\:ds =\sum_{i=1}^{m}\int_{I_i}f\circ\boldsymbol{\sigma}(t)\|\boldsymbol{\sigma}'(t)\|\:dt,  
    \]
    donde $ I_i= (t_i, t_{i+1})$,  $\boldsymbol{\sigma}\lvert_{I_i}$ es de clase $\mathcal{C}^1$  y  $\{t_i\}_{i=0}^{m-1}$  es una partici\'on de $I$.
\end{definition}




\begin{definition}\label{def:int_de_linea}\textbf{Integral de l\'inea para campos escalares.}
    Sea $C\subset \Rn{n}$  una curva simple y suave a trozos  (\ref{param})  y sea $\boldsymbol{\sigma}:I\subseteq\R\to\Rn{n}$ una parametrizaci\'on regular de $C$.   Sea $f:A\subseteq\Rn{n}\to\R$ continua con $C  \subseteq  A$.   Se define la \emph{integral de l\'inea del campo $f$ sobre la curva} $C$, y se nota por $\int_{C}f\:ds$ a
    \[
        \int_{C}f\:ds= \int_{\boldsymbol{\sigma}}f\:ds. 
    \]
   


\end{definition}

\begin{obs} La definici\'on anterior no depende de la parametrizaci\'on.
\end{obs} 

  Veamos un ejemplo para ilustrar la obseravaci\'on  reci\'en enunciada.  

\begin{example}  Sean el campo $f: \mathbb{R}^{2}\to\R$ tal que $f(x,y)=x$ y la curva simple $C$  dada por el segmento que uno los puntos $(0.0)$ y $(1,1)$.  Calcular la integral de l\'inea del campo $f$ sobre la curva $C$.
   
   Por un lado,  es sencillo ver que    $\boldsymbol{\sigma}_1(t)=(t,t)$ con $t\in[0,1]$ es una parametrización suave y regular  de $C$.  Siguiendo  \ref{def:int_de_linea},  la integral de l\'inea de $f$ sobre $C$ es
    \[
        \int_\Gamma f\;ds=\int_0^1f\circ\boldsymbol{\sigma}_1(t)\|\boldsymbol{\sigma}_1'(t)\|\:dt.
    \]
    Como
    \[
        f\circ\boldsymbol{\sigma}_1(t)=f(t,t)=t 
    \]
    y
    \[
        \boldsymbol{\sigma}_1'(t)=(1,1)\implies\|\boldsymbol{\sigma}_1'(t)\|=\sqrt{1^2+1^2}=\sqrt{2};
    \]
    se llega a 
    \[
        \int_C f\;ds=\int_0^1t\sqrt{2}\;dt=\frac{\sqrt{2}}{2}.
    \]

    Por otro lado, tambi\'en es sencillo ver que  $\boldsymbol{\sigma}_2(t)=(t^2,t^2)$ con $t\in[0,1]$ es una parametrización  suave y regular  de $C$.  Siguiendo   \ref{def:int_de_linea},  la integral de l\'inea de $f$ sobre $C$ es
    \[
        \int_C f\;ds=\int_0^1f\circ\boldsymbol{\sigma}_2(t)\|\boldsymbol{\sigma}_2'(t)\|\:dt,
    \]
    Como 
    \[
        f\circ\boldsymbol{\sigma}_2(t)= f(t^{2},t^{2}) = t^2
    \]
    y
    \[
        \boldsymbol{\sigma}_2'(t)=(2t,2t)\implies\|\boldsymbol{\sigma}_2'(t)\|=\sqrt{(2t)^2+(2t)^2}=2t\sqrt{2};
    \]
     se llega a 
    \[
        \int_C f\;ds=\int_0^1t^22t\sqrt{2} \;dt =2\sqrt{2}\int_0^1t^3\;dt=\frac{\sqrt{2}}{2}.
    \]
    
    Que es igual a lo calculado usando a  $\boldsymbol{\sigma}_1$ como parametrizaci\'on de $C.$
  \end{example}  

\begin{obs} 
 Sea $f:A\subseteq\Rn{2} \to\R$   tal que  $f(x,y)>0 \; \forall (x,y) \in A$ y   sea $C  \subseteq  A$  una curva simple en $\Rn{2}$.   La integral de l\'inea de $f$ sobre $C$  se puede interpretar geom\'etricamente como el ``\'area de la valla de base $C$  y altura $f$". 
\end{obs}


\begin{definition}\textbf{Trabajo realizado por una fuerza constante sobre un objeto.}  
Sea un objeto que se desplaza desde un punto $A$ hasta un punto $B$ por el segmento que los une y  bajo la acci\'on de una fuerza \(\mathbf{F}\) constante. 
Se define el     \emph{trabajo realizado por la fuerza sobre el objeto} como:
\[
W = \mathbf{F} \cdot \mathbf{d} = |\mathbf{F}| \, |\mathbf{d}| \, \cos\theta,
\]
donde
\begin{itemize}
    \item \(\mathbf{d} = \mathbf{r}_B - \mathbf{r}_A\) es el vector de desplazamiento del objeto,
    \item \(|\mathbf{F}|\) es la magnitud de la fuerza,
    \item \(|\mathbf{d}|\) es la magnitud del desplazamiento, y
    \item \(\theta\) es el ángulo entre la fuerza y el desplazamiento.
\end{itemize}
\end{definition}


\textbf{Interpretaci\'on:} el trabajo representa la cantidad de energ\'ia transferida al objeto por la fuerza en la direcci\'on del desplazamiento. 
\begin{itemize}
    \item Si \(0 \leq \theta < 90^\circ\), el trabajo es positivo: la fuerza ayuda al movimiento.
    \item Si \(90^\circ < \theta \leq 180^\circ\), el trabajo es negativo: la fuerza se opone al movimiento.
    \item Si \(\theta = 90^\circ\), el trabajo es cero: la fuerza no realiza trabajo sobre el objeto.
\end{itemize}


\begin{definition}\textbf{Integral de trayectoria de campos vectoriales.} Sea $\boldsymbol{\sigma}:[a,b]\to A\subseteq\Rn{n}$ una trayectoria $\mathcal{C}^1$  y sea $F$ un campo vectorial   $\mathbf{F}:A\subseteq\Rn{n}\to\Rn{n}$   tal que este bien definido y sea continuo el campo $\mathbf{F}\circ\boldsymbol{\sigma}$. Se define la \emph{integral de trayectoria del campo vecorial} $\mathbf{F}$ \emph{a lo largo de la trayectoria} $\boldsymbol{\sigma}$, y se nota, $ \int_{\sigma}\mathbf{F}\cdot d\mathbf{s}$, a como
    \[
        \int_{\sigma}\mathbf{F}\cdot d\mathbf{s}=\int_a^b\mathbf{F}\circ\boldsymbol{\sigma}(t)\cdot\boldsymbol{\sigma}'(t)\:dt. 
    \]
\end{definition}

Si $\boldsymbol{\sigma}(t)$ es de clase $\mathcal{C}^1$ a trozos, se define $ \int_{\sigma}\mathbf{F}\cdot d\mathbf{s}$  de la siguiente manera:  

\begin{definition}   Sea $\boldsymbol{\sigma}:I=[a,b]\to\Rn{n}$ una trayectoria de clase $\mathcal{C}^1$ a trozos   y  sea $F$ un campo vectorial tal que este bien definida y sea continua   la funci\'on $F\circ\boldsymbol{\sigma}$.  Se define la \emph{integral de trayectoria del campo vecorial} $\mathbf{F}$ \emph{a lo largo de la trayectoria} $\boldsymbol{\sigma}$, y se  nota $ \int_{\sigma}\mathbf{F}\cdot d\mathbf{s}$, como
  
    \[
       \int_{\sigma}\mathbf{F}\cdot d\mathbf{s}  =\sum_{i=1}^{m}\int_{I_i}\mathbf{F} \circ\boldsymbol{\sigma}(t) \cdot \boldsymbol{\sigma}'(t) \:dt,  
    \]
    donde $ I_i= (t_i, t_{i+1})$,  $\boldsymbol{\sigma}\lvert_{I_i}$ es de clase $\mathcal{C}^1$  y  $\{t_i\}_{i=0}^{m-1}$  es una partici\'on de $I$.
\end{definition}


 \begin{tarea}Sea $\boldsymbol{\sigma}: [-1,1]\to\Rn{2}\;  : \sigma(t) = ( t^{2}, t^{2})$   Calcular, indicando la orientaci\'on elegida para $C$  $$ \int_{C}  (y, 2xy +1)  \cdot d\mathbf{s} $$
 \end{tarea}



\begin{definition}\label{int_cv} \textbf{Integral de l\'inea para campos vectoriales.}
  Sea $C^{+} \subseteq  \Rn{n}$  una curva simple,  suave a trozos  (\ref{param}) orientada (\ref{orien})  y sea $\boldsymbol{\sigma}:I\subseteq\R\to\Rn{n}$ una parametrizaci\'on regular de $C$ que preserva su orientaci\'on  (\ref{orien_pre}). Sea $\mathbf{F}:A\subseteq\Rn{n}\to\Rn{n}$ un campo vectorial continuo sobre $C  \subseteq  A.$  Se define la \emph{integral de} $\mathbf{F}$ \emph{sobre la curva} $\boldsymbol{C}^{+}$ como
    \[
        \int_{C^{+}}\mathbf{F}\cdot d\mathbf{s}=\int_{\boldsymbol{\sigma}}\mathbf{F}\cdot d\mathbf{s} 
    \]
\end{definition}

\begin{obs} Bajo las condiciones de \ref{int_cv}, si $\boldsymbol{\beta}$ es una parametrizaci\'on regular de $C$ que no  preserva (o invierte) su orientaci\'on entonces: 
 \[
        \int_{C^{+}}\mathbf{F}\cdot d\mathbf{s}= - \int_{\boldsymbol{\beta}}\mathbf{F}\cdot d\mathbf{s} 
    \]
\end{obs}







\begin{definition} \textbf{Trabajo de un campo de fuerzas sobre una trayectoria.}
Sea un objeto que se mueve a lo largo de una curva suave y orientada $C$ bajo la acci\'on de un campo de fuerzas $\mathbf{F}$. 
Se define el     \emph{trabajo realizado por el campo sobre el objeto} como la integral de línea:
\[
W = \int_{\sigma}\mathbf{F}\cdot d\mathbf{s}
\]

Si la curva $C$ está parametrizada por $\boldsymbol{\sigma}(t)$ con $t \in [a,b]$, entonces:
\[
W = \int_a^b \mathbf{F}(\boldsymbol{\sigma}(t)) \cdot \boldsymbol{\sigma}(t)'(t) \: dt.
\]
\end{definition}





  \begin{propertie} \textbf{Algebra de Integral de l\'inea.}\label{alge} Sea $C^{+}$ una curva simple, orientada y suave a trozos, adem\'as,  sean  $\mathbf{F}$ y $\mathbf{G}$ campos vectoriales continuos con dominios tales que incluyen a la curva $C$, entonces: 
\begin{itemize}
\item[1.]  \[
        \int_{C^{+}}(\mathbf{F} + \mathbf{G})  \cdot d\mathbf{s}=  \int_{C^{+}} \mathbf{F} \cdot d\mathbf{s}  +  \int_{C^{+}} \mathbf{G}  \cdot d\mathbf{s} 
    \]

\item[2.]  \[
        \int_{C^{+}} k \mathbf{F}  \cdot d\mathbf{s}= k \int_{C^{+}} \mathbf{F} \cdot d\mathbf{s}   
    \]

\item[3.]  \[
        \int_{C^{+}} k \mathbf{F}  \cdot d\mathbf{s}= - \int_{C^{-}} \mathbf{F} \cdot d\mathbf{s}   
    \]

\item[3.]  Si $C = \cup_{i=1}^{n} C_{i},$ entonces \[
        \int_{C^{+}}  \mathbf{F}  \cdot d\mathbf{s}=  \sum_{i=1}^{n}  \int_{C_{i}^{+}} \mathbf{F} \cdot d\mathbf{s}   
    \]

\end{itemize}
 \end{propertie}



 \begin{tarea} Sea $C$ la curva simple dada por la uni\'on de los segmentos que unen los puntos $(0,0)$,  $(0,3)$ y $(0,2)$. Calcular, indicando la orientaci\'on elegida para $C$  $$ \int_{C}  (y^{2}, 2xy +1)  \cdot d\mathbf{s} $$
 \end{tarea}

 \begin{tarea} Mismo ejercicio con la orientaci\'on para $C$ contraria  a la elegida. 
 \end{tarea}

\begin{tarea}Reescribir  (\ref{alge})  para campos escalares. 
\end{tarea}


\begin{definition}
    Si $C$  es una curva simple cerrada,  se suele notar a la integral como
    \[
        \oint_{C}\mathbf{F}\cdot d\mathbf{s},  
    \]  y se la llama integral de circulaci\'on de $\mathbf{F}$ sobre $\Gamma$.
\end{definition}


 \begin{nota}\label{not_int} Sean  $\mathbf{F}=(f_1, f_2)$ y $\mathbf{G}=(g_1, g_2, g_3)$ campos vectoriales, notaremos $$ \int_{C_1}\mathbf{F}\cdot d\mathbf{s} =  \int_{ C_1}f_1\:dx+ f_2 \:dy$$ 
  $$ \int_{C_2}\mathbf{G}\cdot d\mathbf{s} =  \int_{C_2} g_1\:dx+g_2\:dy+g_3\:dz.$$
 \end{nota}
 
 \begin{nota} 
  $\int_{\boldsymbol{\sigma}}\mathbf{F}\cdot d\mathbf{s}=\int_{\boldsymbol{\sigma}}\mathbf{F}\cdot{\mathbf{t}}\:ds$,  donde ${\mathbf{t}}$ es versor tangencial a la curva  $\Gamma = Img(\boldsymbol{\sigma}).$
 \end{nota}
